\section{Numerical Examples}
For all experiments, we use MATLAB on a MacBook Pro with an M4 Max processor with 24GB of RAM, and  YALMIP \cite{lofbergYALMIPToolboxModeling2004} and MOSEK \cite{mosekapsMOSEKOptimizationToolbox2025} to solve SDPs. 
Code is available at {\small \url{https://github.com/naoya-kumagai/sqrt_covariance_steering/}}.
% The SCP parameters are from \cite{oguriSuccessiveConvexificationFeasibility2023a}, summarized in \cref{tab:scp_parameters}. 
For the proposed method, the initial guess \( \bar{\bm{S}} \) is obtained by the covariance matrix interpolation technique from \cite{linRiemannianGeometrySymmetric2019}. The initial guess \( \bar{\bm{L}} \) is set to zero.

% \begin{table}[t]
% 	\centering
% 	\caption{\texttt{SCvx*} parameters ($\epsilon_{\mathrm{opt}}, \epsilon_{\mathrm{feas}}$ defined in text)}
% 	\label{tab:scp_parameters}
% 	\begin{tabular}{cccc}
% 	\toprule
% 	 $\{\rho_0, \rho_1, \rho_2\}$ & $\{\alpha_1, \alpha_2, \beta, \gamma\}$ & $\{r_{\mathrm{init}}, r_{\mathrm{min}}, r_{\mathrm{max}}\}$ & $w_{\mathrm{init}}$\\
% 	\midrule
% 	 $\{0, 0.25, 0.7\}$ & $\{2, 3, 2, 0.9\}$ & $\{0.1, 10^{-6}, 1\}$ & $100$\\
% 	\bottomrule
% 	\end{tabular}
% \end{table}

\subsection{Optimality and Scalability for the Unconstrained Case} \label{sec:comparison_unconstrained}
Here, we compare the horizon-length scalability and solution optimality of our method against Liu et al. \cite{liuOptimalCovarianceSteering2025} and Okamoto and Tsiotras \cite{okamotoOptimalStochasticVehicle2019}. We consider a three-dimensional double integrator with horizon of \(T = 3\) seconds, and vary the horizon length from 10 to 160 (so the step size is \(\Delta t = T / N\)). 
The time-invariant dynamics are given by
\begin{equation} \label{eq:double_integrator_equations}
	A = \smqty[
		I_3 & \Delta t I_3 \\
		0_{3 \times 3} & I_3
	],\ 
	B = \smqty[
		0.5 \Delta t^2 I_3 \\
		\Delta t I_3
	],\ 
	G = \sqrt{q \Delta t} I_6
\end{equation}
with \( q = 0.05 \) being the process noise spectral density.
The initial and terminal state covariance are \(I_6\) and \(0.5 I_6\), respectively. We use an EoQ cost with \(Q = 0.1 I_6\) and \(R = I_3\).
For our method, we set the convergence and feasibility tolerance to \( \epsilon_{\mathrm{opt}} = \epsilon_{\mathrm{feas}} = 10^{-4}\) and scaling matrix \(D_X\) to identity. To time the solutions, we perform 5 trials for each case. The results are shown in \cref{fig:comparison_of_methods}.

For this unconstrained problem, \cite{liuOptimalCovarianceSteering2025} offers fast and optimal solutions. The computational complexity of \cite{okamotoOptimalStochasticVehicle2019} grows quickly, which matches the observations in \cite{rapakouliasDiscreteTimeOptimalCovariance2023}. Due to its specific feedback policy, it also compromises optimality. In comparison, our approach has moderate computational complexity, with rate of growth similar to \cite{liuOptimalCovarianceSteering2025}. The optimality empirically validates \cref{th:global_optimality}. The original cc-CS method \cite{okamotoOptimalCovarianceControl2018} using state history feedback showed much more rapid growth in computational cost and is not included in this comparison.

\begin{figure}[t]
	\centering
	\begin{subfigure}{0.4\linewidth}
		\centering
		\includegraphics[width=\textwidth]{horizon_size_scalability_small.pdf}	
	\end{subfigure}%
	\begin{subfigure}{0.4\linewidth}
		\centering
		\includegraphics[width=\textwidth]{horizon_size_cost_comparision_small.pdf}	
	\end{subfigure}
	\caption{Comparisons for the unconstrained case. (\textbf{Left}) Comparison of the average (in solid lines) and the min-max range (in shaded regions) of the solution times between \cite{liuOptimalCovarianceSteering2025,okamotoOptimalStochasticVehicle2019}, and the proposed method (\textbf{Right}) Comparison of the costs for \cite{okamotoOptimalStochasticVehicle2019} and the proposed method, relative to \cite{liuOptimalCovarianceSteering2025}}
	\label{fig:comparison_of_methods}
\end{figure}

% \begin{remark}
	% We further tested global optimality for unconstrained problems by randomly generating the system matrices. We observed that depending on the dynamics and algorithm hyperparameters, the algorithm may require a large number of iterations (few hundreds) to converge to the global optimum. Furthermore, if using rather large tolerance parameters (e.g. $10^{-4}$), the algorithm may terminate at a suboptimal point; in this case, continuing until convergence with a smaller value (e.g. $10^{-7}$) will more likely achieve global optimality. 
	% \textr{Add some analysis}
% \end{remark}

\subsection{Optimality and Scalability for the Constrained Case}
Next, we perform a comparison of the same methods for path planning in an environment with a circular obstacle. We conservatively approximate the obstacle as a time-varying hyperplane \cite{dicairanoModelPredictiveControl2012}.
These hyperplanes are obtained by solving the deterministic version of the problem to obtain position histories for the linearization reference.
The setting is a 2D double integrator with $q = 0.005$ in \cref{eq:double_integrator_equations}. Between the boundary positions, there is a circular obstacle with its center at $[5, 0]$ and radius $1.2$. We define the nonconvex state CC: $\P{\norm{p - [5, 0]}_2 \geq 1.2} \geq 1 - 0.005$
for each node, where $p = [x, y]^\top$ is the position vector.
The control is chance-constrained as $\P{u \leq 0.15} \geq 1 - 0.005$, $\P{u \geq -0.15} \geq 1 - 0.005$
for each control component and node. The other parameters are $T = 30$, $ \mu_{\mathrm{init}} = [0, 0, 0, 0]^\top$, $ \mu_{\mathrm{fin}} = [10, 0, 0, 0]^\top$, $P_{\mathrm{init}} =\diag{0.1, 0.1, 0.01, 0.01}$, $P_{\mathrm{fin}} = \diag{0.04, 0.04, 0.01, 0.01}$.

When using Liu's formulation \cite{liuOptimalCovarianceSteering2025}, affine CCs are nonconvex in the variables due to the square root term that appears. Two methods exist for handling this nonconvexity: 1) providing a constant reference covariance matrix for all nodes \cite{rapakouliasDiscreteTimeOptimalCovariance2023}, and 2) rewriting the constraint as a difference-of-convex form, solved using the convex-concave procedure (CCP)\cite{pilipovskyComputationallyEfficientChance2024}. We use the latter method, as the former method conservatively approximates the available solution space; this resulted in \textit{artificial infeasibility}. A naive CCP also encounters this issue. Hence, we apply the Penalty CCP method \cite{lippVariationsExtensionConvex2016} which introduces positive slack variables and augments a penalty term in the objective function. The implementation of the Penalty CCP method follows the algorithm of \cite{lippVariationsExtensionConvex2016}. The iterative procedure ends with similar criteria as the proposed method: 1) the nonconvex CC violation and all slack varialbes are smaller than $\epsilon_{\mathrm{feas}}$ and 2) the difference of the augmented objective function with the previous step is smaller than $\epsilon_{\mathrm{opt}}$. The slack variables are multiplied a factor $\tau$ and added to the objective function. $\tau$ is initially 100 and increased by a factor of 2 at each iteration. For this formulation and the proposed method, we set $\epsilon_{\mathrm{opt}} = 10^{-2}$, $\epsilon_{\mathrm{feas}} = 10^{-4}$. 

\begin{figure}[t]
	\centering
	\begin{subfigure}{0.4\linewidth}
		\centering
		\includegraphics[width=\textwidth]{horizon_size_scalability_single_obstacle.pdf}	
	\end{subfigure}%
	\begin{subfigure}{0.4\linewidth}
		\centering
		\includegraphics[width=\textwidth]{horizon_size_cost_comparison_single_obstacle.pdf}	
	\end{subfigure}%
	\caption{Comparisons for the obstacle-avoidance case. (\textbf{Left}) Average (in solid lines) and the min-max range (in shaded regions) of the solution times. (\textbf{Right}) Objective function divided by the horizon length $N$}
	\label{fig:comparison_of_methods_obstacle}
\end{figure} 

\begin{figure}[t]
	\centering
	\begin{subfigure}{0.5\linewidth}
		\centering
		\includegraphics[width=\textwidth]{trajectory_fullcov_N40.pdf}
	\end{subfigure}%
	\begin{subfigure}{0.5\linewidth}
		\centering
		\includegraphics[width=\textwidth]{trajectory_sqrtqr_N40.pdf}
	\end{subfigure}
	\caption{Obstacle avoidance with Liu's formulation with Penalty CCP (left) and proposed method (right) for the case with $N=40$}
	\label{fig:planning_with_obstacles_comparison}
\end{figure}

The comparison results are shown in \cref{fig:comparison_of_methods_obstacle}. 
% The block-matrix method \cite{okamotoOptimalStochasticVehicle2019} is infeasible for all cases, which highlights the limited solution space using its specific feedback control policy. Nevertheless, we include the computation time, which still grows quickly as $N$ increases.
For this constrained problem, the proposed method and CCP-based method have similar computational complexity. The proposed method achieves lower cost. \cite{okamotoOptimalStochasticVehicle2019} reported infeasibility for all cases; since MOSEK's interior-point solver required a typical number of iterations before terminating, we include it as a reference.
\cref{fig:planning_with_obstacles_comparison} compares the trajectory and covariance ellipses for the case with $N=40$. The two methods converge to different solutions, with a difference in how large the position covariance is allowed to grow.

%  The control profiles are shown in \cref{fig:planning_with_obstacles_comparison_control}. A notable difference is seen in the profile of the nominal control magnitude. While the CCP-based method shows a continuous control profile, the proposed method shows large changes in the nominal control magnitude, similar to the deterministic reference case.
\begin{remark}
	The proposed method solves cc-CS with approximately the same computational effort as unconstrained CS.
\end{remark}

\subsection{Numerical Reliability}
Here, we demonstrate the numerical reliability of the proposed method when the variance of some state components are small. We consider a problem similar to \cite{oguriChanceConstrainedControlSafe2024a} which uses CS for a spacecraft rendezvous setting. The problem uses a QoN objective function with $W_k^x = 0$ and $W_k^u = I$ for all $k$ \cref{eq:objective_function} which is common for optimizing the quantile of the spacecraft's fuel consumption. We emphasize that the parameters of uncertainty are rather realistic for a spacecraft rendezvous setting, and not artificial: initial mean position (velocity) $[-3.0, 0.126, 0]^\top$ km ($0_{3 \times 1}$ km/s); initial position (velocity) standard deviation $100$ m ($1.0$ m/s); terminal mean position (velocity) $[0.0, 0.05, 0]^\top$ km ($0_{3 \times 1}$ km/s); terminal position(velocity) standard deviation $10.0$ m ($0.1$ m/s); stochastic acceleration $1.0$ mm/s$^{3/2}$; horizon length $N = 14$; time step $30$ s; chief orbit radius $7228$ km; gravitational parameter $3.986 \times 10^{5}$ km$^3$/s$^2$.

The continuous-time Clohessy-Wiltshire-Hill equations \cite{schaubAnalyticalMechanicsSpace2003} are discretized using zero-order hold control, as described in \cite{oguriChanceConstrainedControlSafe2024a}. 
We use $D_X = \mathrm{diag}(10, 10, 10, 10^3, 10^3, 10^3)$ for the trust region scaling, chosen to make the magnitude of $D_X S_k$ close to unity. We set \( \epsilon_{\mathrm{opt}} = \epsilon_{\mathrm{feas}} = 10^{-5}\).
The proposed method solves the problem in 42 iterations in 5.0 seconds (average of 10 runs). The results are shown in \cref{fig:example_cwh_trajectory}, along with the results of 1000 Monte Carlo simulations. We can observe that the terminal covariance constraint is satisfied with equality, and the Monte Carlo samples are within the 3-sigma bounds.

\begin{figure}[t]
	\centering
	\includegraphics[width=\linewidth]{example_cwh_trajectory.pdf}
	\caption{Projection of trajectory onto the $x$-$y$ plane for the spacecraft rendezvous problem, with $3\sigma$ covariance ellipses around the mean trajectory}
	\label{fig:example_cwh_trajectory}
\end{figure}
\begin{figure}[t]
	\centering
	\includegraphics[width=0.75\linewidth]{example_cwh_covariance_propagation_loss.pdf}
	\caption{Covariance propagation loss for the spacecraft rendezvous problem}
	\label{fig:example_cwh_covariance_propagation_loss}
\end{figure}


For the approach of \cite{liuOptimalCovarianceSteering2025}, we use $\eta = 10^{-3}$ for the convexifying term (see \cref{remark:convexifying_term}).
For this scenario, MOSEK returns status \texttt{UNKNOWN}; this indicates numerical issues, where the returned solution may or may not be close to a feasible and optimal solution \cite{mosekapsMOSEKOptimizationToolbox2025}. 
Here, the covariance dynamics are largely violated, according to a measure of loss in covariance propagation defined as follows:
\begin{equation}
	\text{Loss}_k = \norm{\phi(K_k, P_k) - P_{k+1}}_F / \norm{P_{k+1}}_F
\end{equation}
where \( K_k, P_k \) are the solution of the optimization, and $\phi(K_k, P_k)$ is the RHS of \cref{eq:full_covariance_propagation}. 
\cref{fig:example_cwh_covariance_propagation_loss} shows the loss for the two methods, showing the superior accuracy of the proposed method's covariance propagation; using tighter convergence tolerance for the proposed method can further improve the accuracy if desired. We also tested with different values of $\eta$ and alternative objective functions for the full-covariance approach, such as the EoQ form, \( \sum_k \norm{v_k}_2 + \lambda_{\max}(Y_k) + \eta \trace(Y_k)\), and \(\sum_k \norm{v_k}_2 + \trace(Y_k) \), all either resulting in numerical issues and/or lossy covariance propagation.
These results support our hypothesis: by working with the square root of the covariance matrix, solvers can work with better-scaled problems and are less likely to encounter numerical issues.